\section{Spin precession in the axion search mode}
\label{sec:axion}

Peccei--Quinn axions, proposed as a solution of the strong CP problem, are
among the most likely dark matter candidates. A storage ring can be used as an
antenna for the search for the galactic axion~\cite{Nikolaev:2022wyi}.

Current searches for resonant spin rotation in storage rings use the method
developed by the JEDI collaboration, in which the growth of the vertical
polarization out of the horizontal plane indicates the presence of the axion
field~\cite{Nikolaev:2022wyi}. Based on an analytical treatment of the effect of the spin coherence
time on the frequency scan for the axion signal, an alternative scheme has been
proposed in which an initially vertical spin is rotated into the horizontal
plane~\cite{Nikolaev:2022wyi}. This scheme is free from the ambiguity of the
axion field phase, requires no RF spin flippers, and can be readily
implemented with both deuterons and protons stored in the Nuclotron and NICA
rings.

Consider the proposed experiment on the example of the deuteron EDM search
with polarized beams. It is based on bypasses with Wien filters installed on
them (Section~\ref{sec:nica-bypass}); the QFS lattice of NICA with bypasses is
therefore of particular interest. The EDM and axion experiments differ in that
the axion search does not require the frozen spin condition.

The method relies on the interaction of the axion with the EDM of the particles
in the bunch. This interaction modifies the spin precession described by the
T-BMT equation~\eqref{eq:tbmt} and increases the vertical spin component. Since
the axion mass is unknown, it has to be found by scanning the spin precession
frequency, which is done by varying the fields of the Wien filters.

In the bending arcs of the NICA collider, the particle momentum rotates by
$\Phi^B_{arc} = \pi$, while the deuteron spin ($\gamma G<0$) rotates in the
opposite direction in the horizontal plane relative to the momentum by the
angle $\Phi_s^{arc}$ of Eq.~\eqref{eq:phi_arc} (Fig.~\ref{fig:axion_diagram}).
In each of the two straight sections, the Wien filters rotate the spin back as
described by Eqs.~\eqref{eq:phi_e}--\eqref{eq:wien}. For a maximum electric
field of \SI{120}{kV/cm}, the required magnetic field is below \SI{80}{mT}.

\begin{figure}[htbp]
  \centering
  \includegraphics[width=0.75\textwidth]{axion_spin_diagram}
  \caption{Vector diagram of the spin transformation for the deuteron.}
  \label{fig:axion_diagram}
\end{figure}

In the axion search mode, the QFS condition~\eqref{eq:qfs} is modified to
\begin{equation}
  (\gamma G+1)\,\Phi^B_{ss} - \left(\gamma G+\frac{\gamma}{\gamma+1}\right)\beta^2\,\Phi^E_{ss} = \Delta\Phi + \pi\gamma G ,
  \label{eq:axion_qfs}
\end{equation}
where $\Delta\Phi$ is the detuning phase with respect to the QFS case
$\Delta\Phi = 0$. By changing the energy and the Wien filter fields $B_{ss}$
and $E_{ss}$, one changes the frequency of the spin precession around the
vertical axis,
\begin{equation}
  f_{spin} = \frac{\Delta\Phi\,c\beta}{C_{ring}},
  \label{eq:axion_fspin}
\end{equation}
where $C_{ring}$ is the ring circumference, $c\beta$ is the beam velocity and
$\Delta\Phi$ is the detuning phase defined by Eq.~\eqref{eq:axion_qfs}. For
example, at $\Delta\Phi = \SI{1}{rad}$ the spin precession frequency around the
vertical axis is \SI{270}{kHz}, i.e., the axion search range is
$f_{spin} = 0$--\SI{270}{kHz}.\todo[inline]{Eq.~\eqref{eq:axion_fspin} gives an
angular frequency (rad/s); the frequency in Hz is $f_{spin}/2\pi$.}

For a proton beam, the QFS option cannot be realized in NICA: the proton
magnetic anomaly $G_p = 1.793$ is an order of magnitude larger, and the spin
rotates by more than \ang{360} in each arc, so that the accumulated EDM signal
practically vanishes. For the axion search, however, the idea remains valid,
in a different range of $f_{spin}$.

The material of this section has been reported
in~\cite{Nikolaev:2022wyi,Senichev:2022ide,Aksentev:2026mjw}.
